% Compile with XeLaTeX (Overleaf: Menu > Compiler > XeLaTeX) for the Vietnamese name.
\documentclass[11pt,a4paper]{article}
\usepackage[margin=2cm]{geometry}
\usepackage{fontspec}
\usepackage{array,tabularx,booktabs}
\usepackage{longtable}
\usepackage[hidelinks]{hyperref}
\usepackage{xurl}
\setlength{\parindent}{0pt}
\setlength{\parskip}{5pt}
\usepackage{ragged2e}
\pagestyle{empty}

\newcommand{\sect}[2]{\textbf{#1} \textbf{(#2 points)}\par}
\newcommand{\hint}[1]{{\small\itshape #1}\par}
\newcommand{\X}{$\times$}

\begin{document}

\begin{center}
{\LARGE\textbf{MIDTERM REPORT}}\\[4pt]
{\large\textbf{Undergraduate Research Opportunities Program 1 (UROP1010)}}
\end{center}

\begingroup\small\renewcommand{\arraystretch}{1.25}\raggedright
\begin{tabular}{|>{\bfseries\raggedright\arraybackslash}p{3.7cm}|>{\raggedright\arraybackslash}p{5.3cm}|>{\bfseries\raggedright\arraybackslash}p{1.8cm}|>{\raggedright\arraybackslash}p{4.3cm}|}
\hline
Student name: & Đặng Tuấn Kiệt & Date: & 13/11/2026 \\ \hline
Student ID: & V202502041 & Year: & 2 \\ \hline
College: & College of Engineering and Computer Science & Program: & Bachelor of Data Science \\ \hline
Faculty supervisor name: & \multicolumn{3}{>{\raggedright\arraybackslash}p{12.1cm}|}{Dr. Nguyễn Hà Thành} \\ \hline
Research Project Title: & \multicolumn{3}{>{\raggedright\arraybackslash}p{12.1cm}|}{Automated Counterfactual Challenge Generation for Legal QA Systems (working title)} \\ \hline
\end{tabular}
\endgroup

\bigskip
%------------------------------------------------------------
\sect{I. Project Overview}{20}
\hint{Please provide a brief overview of the research project, outlining its goals, objectives, and significance in the respective field of study.}

Large language models are increasingly used for legal question answering and often give fluent explanations alongside their answers. Such explanations can sound reasonable without showing how the answer depends on the facts of a case. This project explores a behavior-based alternative to judging explanations by plausibility: automatically generating counterfactual factual challenges, meaning legally meaningful changes to the facts of a case, and observing how a Legal QA system's conclusion and explanation respond.

The working direction considers three evaluation views: (1) relevant-fact sensitivity, i.e.\ whether the conclusion changes appropriately when a legally relevant fact is changed; (2) irrelevant-fact robustness, i.e.\ whether it stays stable when an irrelevant fact is added; and (3) explanation alignment, i.e.\ whether the explanation is consistent with the observed change in the decision. These are research directions, not yet validated metrics.

The research question is not yet finalized. Three candidates (evaluation-focused, challenger-focused, combined) are under discussion with the supervisor, with a planned scope of 3--4 months. The goal of this first phase is to narrow the question to one that is precise and experimentally testable, through a focused reading of four papers recommended by the supervisor. This is a preliminary research direction: no Challenger system has been built, and no novelty claim is made.

\bigskip
%------------------------------------------------------------
\sect{II. Progress Update}{50}
\hint{Please write an overview of your tasks and progress in the research project, detailing your accomplishments, challenges faced, and the steps taken to overcome these challenges. Include any significant findings or developments in the project.}

\textbf{Task 1: } Initial literature review on explanation faithfulness and counterfactual evaluation. I read the survey on faithful model explanation in NLP (\textit{Computational Linguistics}, 2024) and the ACL 2024 paper on faithfulness versus self-consistency of natural language explanations.

Findings: an explanation can be plausible without being faithful; counterfactual interventions are attractive but have known limits (out-of-distribution inputs, interacting features, hindsight bias, no agreed faithfulness metric); and tests that compare an answer with its explanation under input edits measure self-consistency, not faithfulness. As a result, the project reports observable decision behavior and behavioral consistency rather than faithfulness.

\textbf{Task 2: } Research direction refinement based on supervisor feedback. The question moved from a broad one about legal explanation faithfulness to the current direction: generating legally meaningful factual interventions and observing how the system responds.

I defined relevant versus irrelevant challenges and the three evaluation views above, and recorded five open methodological problems: defining a ``legally relevant fact'', establishing the expected outcome, evaluating the Challenger itself, choosing intervention strength, and handling interacting legal factors. Three candidate research questions were written down; none is selected yet.

\textbf{Task 3: } Investigation of the four recommended papers and methodological questions. I read all four papers in full and wrote structured notes (problem, terminology, method, assumptions, limitations, relevance to the project), separating what each paper reports from my interpretation and from proposals.

The ICLR 2025 paper on faithfulness of LLM explanations offers a pipeline of concept extraction, counterfactual generation (replace or remove a concept) and quality checks (coherency, completeness, minimality). It does not test legal data and only edits existing concepts rather than adding new legally meaningful facts.

The legal reasoning paper (\textit{Computer Law \& Security Review}, 2025) uses a fact / theory / hypothesis / outcome framing in which hypothetical facts challenge a legal theory, and points to factors and dimensions from case-based reasoning as a possible way to operationalize legal relevance. Its own evaluation uses LLM judges without reference answers on a small number of prompts, so I do not adopt it as an evaluation model.

None of the four papers gives a procedure for establishing the expected legal outcome after a fact is changed; this remains the main open problem. Whether it is a research gap has not been established, since literature beyond the four papers has not been reviewed. Challenge: the full text of the legal reasoning paper was blocked for automated access; I resolved it by downloading the open-access PDF in a normal browser. The research question and the task definition still require discussion with the supervisor.

\textbf{Code reference:} \url{https://github.com/Cuki2910/legal_cf_bench/tree/report-v2-validation} (public; v1/v2 code and aggregates pushed on branch report-v2-validation).

\textbf{Task 4: } Preliminary synthetic-rule pilot and benchmark validation. I implemented a template-based prototype that supplies a rule and case facts, generates controlled factual interventions, and records model answers. It uses a three-valued evaluator (yes / no / undetermined), with missing facts treated as unknown rather than false. These templates are fixed interventions, not an automated LLM-based Challenger.

I completed an exploratory evaluation of three model routes under two prompt wordings on a held-out-at-run-time set of eight synthetic schemas (239 items per prompt; 1,434 responses in total). A subsequent audit found that several factual descriptions did not determine the atom values assumed by the symbolic gold, and that multiple item IDs shared identical prompts. For example, depositing a watch in a hotel safe on the first day does not establish that it remained there when it was stolen. Some answers marked wrong were therefore defensible as undetermined. I do not interpret these scores as a ranking of model ability or as evidence of legal explanation faithfulness.

The main finding is methodological: computing gold from a formula checks symbolic consistency, but does not validate the mapping from natural-language facts to that formula. I preserved the original data and outputs, documented the ambiguities, and specified repairs: explicit factual conditions, single-proposition interventions, separate stimulus-level and pair-level measures, and human review before another evaluation. The revised offline implementation is complete and passes regression tests; its natural-language gold and model performance have not yet been validated. Detailed numerical tables, uncertainty estimates, archived-run inventory and the repository link are included in Appendix A. The next step is to agree on a narrower research question with the supervisor and complete the data-validation protocol before further model evaluation.

\bigskip
%------------------------------------------------------------
\sect{III. Self-Evaluation}{30}
\hint{Please assess your performance during the project based on the following evaluation criteria on a scale of 1 to 5 and make specific comments in each category.\\
(1-Unsatisfactory performance; 2-Below expectation; 3-Average performance; 4-Good performance; 5-Excellent performance)}

\small
\renewcommand{\arraystretch}{1.3}
\begin{tabularx}{\textwidth}{|>{\raggedright\arraybackslash}p{4.6cm}|*{5}{>{\centering\arraybackslash}p{0.55cm}|}X|}
\hline
\textbf{Criteria} & \textbf{1} & \textbf{2} & \textbf{3} & \textbf{4} & \textbf{5} & \textbf{Comment} \\ \hline
1 -- I have maintained professionalism during the program (professional presence, punctuality, etc.) & & & & \X & & Attended the weekly reading seminar almost every week and kept written notes of decisions and readings. \\ \hline
2 -- I have been able to complete given tasks with timeliness, accuracy, quality, and efficiency. & & & & \X & & Completed the assigned work for this stage: read all four recommended papers in full, wrote structured notes with sources, and drafted candidate research questions. Progress was slower than ideal because I am also working on other research projects, and the research question is not yet finalized. I plan to dedicate more time in the second half. \\ \hline
3 -- I have demonstrated analytical thinking, critical thinking, and problem-solving skills while completing given tasks. & & & & & \X & Separated what each paper reports from my interpretation, and flagged unverified claims instead of guessing. \\ \hline
4 -- I have proposed new ideas, creative initiatives and/or applied innovative methods/ways to complete given tasks. & & & & \X & & Proposed the relevant/irrelevant challenge framing and a link between legal factors and hypothetical facts. I implemented a preliminary synthetic-rule pilot to probe this framing; its audit exposed data-validity issues, so it does not yet validate the proposed evaluation. \\ \hline
5 -- I have demonstrated effective communication and collaboration with peers and supervisors (listen sincerely, be willing to help, be trusted, etc.) & & & \X & & & Incorporated the supervisor's feedback into the research direction and attended the weekly reading seminar almost every week. However, I did not communicate often enough with my supervisor outside the seminar, partly because I was busy with other research projects. I take responsibility for this and will update my supervisor more regularly and prepare questions before each discussion. \\ \hline
\end{tabularx}
\normalsize

\medskip
\textbf{What are your main strengths?}\\
I have hands-on research experience from several other projects at the VIN-SSUM lab on sustainable mobility, including deep-learning modelling of motorcycle risky behavior, a mobility capability study with field data collection, a world-model feasibility study, and social-media data crawling pipelines. This experience helps me read research papers quickly, extract methods and limitations, and turn them into a focused research direction. I am careful about sources: every figure in my notes is tied to a section or marked as unverified, and I keep claims within what the evidence supports. I read papers critically, separate findings from interpretation, and found another route when access to a paper was blocked. I also took part in the weekly reading seminar almost every week.

\textbf{What are your main areas of improvement?}\\
Time management is my main area of improvement: I am also involved in other research projects at the VIN-SSUM lab, so I could not give this project my full effort in the first half. I will schedule fixed weekly time for it and commit to more consistent progress in the second half. The research question is still open, so the next step is to narrow it with my supervisor and define the Challenger task and evaluation protocol. I also need to read beyond the four papers (for example case-based reasoning work such as HYPO and CATO) before judging any research gap, and to turn the exploratory pilot into a carefully validated experiment rather than treating preliminary scores as conclusions.

% BEGIN GENERATED RESULTS
\clearpage

\pagestyle{plain}

\appendix

\section*{Appendix A. Experimental results and implementation evidence}

\addcontentsline{toc}{section}{Experimental results}

\textbf{Interpretation boundary.} All numerical results below are agreement with the current synthetic gold, not validated legal accuracy or a model ranking. Gold was computed from symbolic rules, but a post-run audit found mismatches between some English facts and their intended atom values. Scores are retained without removing disputed items. Human/expert validation of the new item set remains incomplete.

\textbf{Repository.} \url{https://github.com/Cuki2910/legal_cf_bench/tree/report-v2-validation}. The GitHub repository is public. V1/v2 code and aggregate result files have been pushed to branch report-v2-validation. Raw API responses remain local, so the public aggregates do not independently reproduce provider outputs.

\textbf{Local evidence.} Counts were recomputed from \path{legal_cf_bench/results/*/raw.jsonl}; v0 values use the saved \path{summary.json}. The frozen test hash matches \path{docs/TEST_FREEZE.txt}. The accompanying \path{docs/REPORT_RESULTS.csv} and \path{docs/REPORT_RESULTS.json} contain the tabulated aggregates; \path{build_report_results.py} rebuilds this appendix without API calls. No raw response was discarded or retried because it was wrong.

\subsection*{A.1 Setup, data sizes and operational definitions}

The evaluated gateway route labels were \path{cx/gpt-6.1-sol} (Sol), \path{cx/gpt-6-luna[1m]} (Luna), and \path{cc/claude-haiku-4-5-20251001} (Haiku 4.5). These are requested route names; the actual backing model identity was not independently verified. V1 used temperature 0 and a 300-token output limit. Prompt pA and prompt pB differ in wording/order; they are not repeated runs of the same configuration.

\begingroup\small\setlength{\tabcolsep}{4pt}
\begin{longtable}{lrrrr}
\caption{Experiment inventory. Item counts are request-level IDs, not independent observations.}\\
\toprule
Stage & Schemas & Bases & Items/prompt & Saved outputs \\
\midrule\endfirsthead
\toprule
Stage & Schemas & Bases & Items/prompt & Saved outputs \\
\midrule\endhead
\bottomrule\endfoot
v0, 2 models, p1 & 10 & 10 & 30 & 60 \\
v1 dev, 3 models, pA & 7 & 21 & 235 & 705 \\
v1 test, 3 models, pA/pB & 8 & 24 & 239 & 1434 \\
\end{longtable}\endgroup


The 705 dev and 1,434 test outputs refer to the latest saved run directories. Historical development files overlap later retained records and must not be summed as independent API calls. Dev templates were revised after observing outputs; dev scores are diagnostic only. V1 contains 474 item IDs per prompt across dev and test. Test bases are label-balanced (8 yes, 8 no, 8 undetermined).

\begingroup\small\setlength{\tabcolsep}{4pt}
\begin{longtable}{lrr}
\caption{Frozen v1 test composition.}\\
\toprule
Condition & Count & Expected relation to base \\
\midrule\endfirsthead
\toprule
Condition & Count & Expected relation to base \\
\midrule\endhead
\bottomrule\endfoot
Base & 24 & -- \\
F: plain atom flip & 32 & Gold changes \\
N3: alternate-text flip & 32 & Gold changes \\
S: supply missing atom & 16 & U to yes/no \\
N: same-value near variant & 39 & Gold unchanged \\
I1: plausible irrelevant & 48 & Gold unchanged \\
I2: harmless irrelevant & 48 & Gold unchanged \\
\end{longtable}\endgroup


\textbf{Metrics.} M0: base accuracy. M1: accuracy after F/N3/S; M2: accuracy after N; M3a/M3b: accuracy after I1/I2. Unconditional values compare each target answer with gold; they do not by themselves require a change/stay relative to the model's own base answer. Conditional values restrict to correctly answered bases. M4 requires the base and every intervention within a case to be correct. M5 is a limited proxy: on correctly answered flips with determinate gold, does the one reported fact ID equal the edited fact ID? M6 measures answer agreement between pA and pB, not correctness. M7 separates M1 by base-gold direction. Headline metrics were M1, M2 and M3a; other analyses are supplementary.

\textbf{Uncertainty.} V1 intervals use percentile cluster bootstrap over schemas (1,000 resamples, seed 0); counts are shown to expose denominators. Eight test clusters are too few for broad generalization. Paired comparison intervals below use 2,000 schema resamples, keep both prompts in each cluster, and have no multiple-comparison correction. These intervals do not account for disputed gold. Degenerate [100,100] conditional intervals only reflect the observed all-correct subset, not certainty about population performance.

\subsection*{A.2 Original pilot v0 (historical; invalid for performance claims)}

\begingroup\small\setlength{\tabcolsep}{4pt}
\begin{longtable}{lrrrr}
\caption{Saved v0 results; sensitivity/robustness conditioned on base correctness.}\\
\toprule
Model & Base & Sensitivity & Robustness & Keyword proxy \\
\midrule\endfirsthead
\toprule
Model & Base & Sensitivity & Robustness & Keyword proxy \\
\midrule\endhead
\bottomrule\endfoot
Sol & 9/10 (90.0\%) & 7/9 (77.8\%) & 9/9 (100.0\%) & 7/7 (100.0\%) \\
Luna & 10/10 (100.0\%) & 10/10 (100.0\%) & 10/10 (100.0\%) & 10/10 (100.0\%) \\
\end{longtable}\endgroup


V0 had only 10 cases, binary answers, imbalanced base labels, no near-miss category, and ambiguous/manual gold (including a necessary-condition versus sufficient-condition error). Its lexical explanation proxy is not explanation alignment. The results motivated v1 rather than supporting a performance claim. For example, a 10/10 v0 Wilson interval is approximately [72.2,100]\%, illustrating the small sample size.

\subsection*{A.3 Latest development-set results (not held out)}

\begingroup\small\setlength{\tabcolsep}{4pt}
\begin{longtable}{lrrrrr}
\caption{Latest v1 dev scores, seven schemas, prompt pA. Gold/template revisions were informed by earlier outputs.}\\
\toprule
Model & M0 & M1 & M2 & M3a & M3b \\
\midrule\endfirsthead
\toprule
Model & M0 & M1 & M2 & M3a & M3b \\
\midrule\endhead
\bottomrule\endfoot
Sol & 21/21 (100.0\%) & 85/86 (98.8\%) & 42/44 (95.5\%) & 42/42 (100.0\%) & 42/42 (100.0\%) \\
Luna & 20/21 (95.2\%) & 83/86 (96.5\%) & 41/44 (93.2\%) & 42/42 (100.0\%) & 41/42 (97.6\%) \\
Haiku 4.5 & 18/21 (85.7\%) & 79/86 (91.9\%) & 37/44 (84.1\%) & 38/42 (90.5\%) & 38/42 (90.5\%) \\
\end{longtable}\endgroup


\begingroup\small\setlength{\tabcolsep}{4pt}
\begin{longtable}{lrrr}
\caption{Development secondary metrics after the final-answer parser was fixed before test.}\\
\toprule
Model & M4 & M5 proxy & Invalid \\
\midrule\endfirsthead
\toprule
Model & M4 & M5 proxy & Invalid \\
\midrule\endhead
\bottomrule\endfoot
Sol & 18/21 (85.7\%) & 77/79 (97.5\%) & 0 \\
Luna & 13/21 (61.9\%) & 63/77 (81.8\%) & 0 \\
Haiku 4.5 & 9/21 (42.9\%) & 64/77 (83.1\%) & 0 \\
\end{longtable}\endgroup


The early four-schema dev set showed ceiling behavior for Sol; three harder schemas were added. Ambiguous examples were subsequently clarified. Those iterations are development decisions, not additional held-out evidence. Haiku's three earlier ``unparsable'' dev responses contained self-corrections with multiple JSON objects; the parser was fixed to take the last parseable object before freezing test. Final dev values above are recomputed from stored raw with that rule.

\subsection*{A.4 Frozen test: all three models and both prompts}

\begingroup\small\setlength{\tabcolsep}{4pt}
\begin{longtable}{lrrrrr}
\caption{Test results, pA, unconditional. All 239 items retained per model.}\\
\toprule
Model & M0 & M1 & M2 & M3a & M3b \\
\midrule\endfirsthead
\toprule
Model & M0 & M1 & M2 & M3a & M3b \\
\midrule\endhead
\bottomrule\endfoot
Sol & 17/24 (70.8\%) & 61/80 (76.2\%) & 27/39 (69.2\%) & 38/48 (79.2\%) & 38/48 (79.2\%) \\
Luna & 20/24 (83.3\%) & 65/80 (81.2\%) & 31/39 (79.5\%) & 40/48 (83.3\%) & 42/48 (87.5\%) \\
Haiku 4.5 & 21/24 (87.5\%) & 70/80 (87.5\%) & 32/39 (82.1\%) & 43/48 (89.6\%) & 43/48 (89.6\%) \\
\end{longtable}\endgroup


\begingroup\small\setlength{\tabcolsep}{4pt}
\begin{longtable}{lrrrrr}
\caption{Test results, pB, unconditional. All 239 items retained per model.}\\
\toprule
Model & M0 & M1 & M2 & M3a & M3b \\
\midrule\endfirsthead
\toprule
Model & M0 & M1 & M2 & M3a & M3b \\
\midrule\endhead
\bottomrule\endfoot
Sol & 19/24 (79.2\%) & 64/80 (80.0\%) & 27/39 (69.2\%) & 39/48 (81.2\%) & 39/48 (81.2\%) \\
Luna & 20/24 (83.3\%) & 65/80 (81.2\%) & 29/39 (74.4\%) & 40/48 (83.3\%) & 40/48 (83.3\%) \\
Haiku 4.5 & 23/24 (95.8\%) & 69/80 (86.2\%) & 35/39 (89.7\%) & 44/48 (91.7\%) & 45/48 (93.8\%) \\
\end{longtable}\endgroup


\begingroup\small\setlength{\tabcolsep}{4pt}
\begin{longtable}{llrrrrr}
\caption{Schema-cluster 95\% intervals for unconditional test rates (percentage points).}\\
\toprule
Model & Prompt & M0 & M1 & M2 & M3a & M3b \\
\midrule\endfirsthead
\toprule
Model & Prompt & M0 & M1 & M2 & M3a & M3b \\
\midrule\endhead
\bottomrule\endfoot
Sol & pA & [58.3, 83.3] & [63.2, 87.5] & [48.8, 88.9] & [62.5, 91.7] & [62.5, 91.7] \\
Luna & pA & [70.8, 95.8] & [69.0, 94.4] & [66.0, 94.1] & [70.8, 93.8] & [77.1, 95.8] \\
Haiku 4.5 & pA & [70.8, 100.0] & [75.0, 98.8] & [67.9, 97.1] & [75.0, 100.0] & [77.1, 97.9] \\
Sol & pB & [62.5, 91.7] & [67.7, 92.2] & [48.8, 88.9] & [68.8, 93.8] & [68.8, 93.8] \\
Luna & pB & [70.8, 95.8] & [70.7, 92.2] & [65.9, 85.0] & [70.8, 95.8] & [70.8, 95.8] \\
Haiku 4.5 & pB & [87.5, 100.0] & [79.5, 94.1] & [82.6, 97.3] & [79.2, 100.0] & [85.4, 100.0] \\
\end{longtable}\endgroup


\begingroup\small\setlength{\tabcolsep}{4pt}
\begin{longtable}{llrrrr}
\caption{Test metrics conditional on a correct base. Denominators differ by model.}\\
\toprule
Model & Prompt & M1c & M2c & M3ac & M3bc \\
\midrule\endfirsthead
\toprule
Model & Prompt & M1c & M2c & M3ac & M3bc \\
\midrule\endhead
\bottomrule\endfoot
Sol & pA & 38/52 (73.1\%) & 23/23 (100.0\%) & 34/34 (100.0\%) & 34/34 (100.0\%) \\
Luna & pA & 51/62 (82.3\%) & 28/29 (96.6\%) & 38/40 (95.0\%) & 40/40 (100.0\%) \\
Haiku 4.5 & pA & 63/70 (90.0\%) & 29/32 (90.6\%) & 42/42 (100.0\%) & 41/42 (97.6\%) \\
Sol & pB & 47/58 (81.0\%) & 25/26 (96.2\%) & 38/38 (100.0\%) & 38/38 (100.0\%) \\
Luna & pB & 49/62 (79.0\%) & 27/29 (93.1\%) & 40/40 (100.0\%) & 40/40 (100.0\%) \\
Haiku 4.5 & pB & 65/76 (85.5\%) & 34/36 (94.4\%) & 44/46 (95.7\%) & 45/46 (97.8\%) \\
\end{longtable}\endgroup


\begingroup\small\setlength{\tabcolsep}{4pt}
\begin{longtable}{lllrr}
\caption{Conditional test intervals. Selection on correct bases prevents a like-for-like model comparison.}\\
\toprule
Model & Prompt & Metric & Correct/eligible & 95\% CI \\
\midrule\endfirsthead
\toprule
Model & Prompt & Metric & Correct/eligible & 95\% CI \\
\midrule\endhead
\bottomrule\endfoot
Sol & pA & M1 & 38/52 (73.1\%) & [55.0, 87.0] \\
Sol & pA & M2 & 23/23 (100.0\%) & [100.0, 100.0] \\
Sol & pA & M3a & 34/34 (100.0\%) & [100.0, 100.0] \\
Sol & pA & M3b & 34/34 (100.0\%) & [100.0, 100.0] \\
Luna & pA & M1 & 51/62 (82.3\%) & [69.0, 94.4] \\
Luna & pA & M2 & 28/29 (96.6\%) & [90.3, 100.0] \\
Luna & pA & M3a & 38/40 (95.0\%) & [87.5, 100.0] \\
Luna & pA & M3b & 40/40 (100.0\%) & [100.0, 100.0] \\
Haiku 4.5 & pA & M1 & 63/70 (90.0\%) & [77.8, 98.8] \\
Haiku 4.5 & pA & M2 & 29/32 (90.6\%) & [80.6, 100.0] \\
Haiku 4.5 & pA & M3a & 42/42 (100.0\%) & [100.0, 100.0] \\
Haiku 4.5 & pA & M3b & 41/42 (97.6\%) & [92.5, 100.0] \\
Sol & pB & M1 & 47/58 (81.0\%) & [66.7, 92.9] \\
Sol & pB & M2 & 25/26 (96.2\%) & [85.0, 100.0] \\
Sol & pB & M3a & 38/38 (100.0\%) & [100.0, 100.0] \\
Sol & pB & M3b & 38/38 (100.0\%) & [100.0, 100.0] \\
Luna & pB & M1 & 49/62 (79.0\%) & [64.3, 92.2] \\
Luna & pB & M2 & 27/29 (93.1\%) & [82.8, 100.0] \\
Luna & pB & M3a & 40/40 (100.0\%) & [100.0, 100.0] \\
Luna & pB & M3b & 40/40 (100.0\%) & [100.0, 100.0] \\
Haiku 4.5 & pB & M1 & 65/76 (85.5\%) & [77.6, 94.1] \\
Haiku 4.5 & pB & M2 & 34/36 (94.4\%) & [88.9, 100.0] \\
Haiku 4.5 & pB & M3a & 44/46 (95.7\%) & [90.5, 100.0] \\
Haiku 4.5 & pB & M3b & 45/46 (97.8\%) & [93.5, 100.0] \\
\end{longtable}\endgroup


\begingroup\small\setlength{\tabcolsep}{4pt}
\begin{longtable}{llrrrr}
\caption{Secondary test metrics. M5 is reported for completeness, not as validated alignment.}\\
\toprule
Model & Prompt & M4 & M4 CI & M5 proxy & Invalid \\
\midrule\endfirsthead
\toprule
Model & Prompt & M4 & M4 CI & M5 proxy & Invalid \\
\midrule\endhead
\bottomrule\endfoot
Sol & pA & 8/24 (33.3\%) & [8.3, 62.5] & 58/59 (98.3\%) & 0 \\
Luna & pA & 12/24 (50.0\%) & [20.8, 79.2] & 58/63 (92.1\%) & 0 \\
Haiku 4.5 & pA & 15/24 (62.5\%) & [37.5, 83.3] & 58/68 (85.3\%) & 1 \\
Sol & pB & 11/24 (45.8\%) & [16.7, 75.0] & 59/62 (95.2\%) & 0 \\
Luna & pB & 10/24 (41.7\%) & [16.7, 66.7] & 50/63 (79.4\%) & 0 \\
Haiku 4.5 & pB & 14/24 (58.3\%) & [33.3, 83.3] & 58/67 (86.6\%) & 0 \\
\end{longtable}\endgroup


\begingroup\small\setlength{\tabcolsep}{4pt}
\begin{longtable}{lrr}
\caption{M6: prompt stability across the same 239 test IDs.}\\
\toprule
Model & Same answer pA/pB & 95\% CI \\
\midrule\endfirsthead
\toprule
Model & Same answer pA/pB & 95\% CI \\
\midrule\endhead
\bottomrule\endfoot
Sol & 232/239 (97.1\%) & [93.9, 99.2] \\
Luna & 221/239 (92.5\%) & [87.3, 97.6] \\
Haiku 4.5 & 208/239 (87.0\%) & [79.0, 93.8] \\
\end{longtable}\endgroup


Haiku pA has one invalid output at \path{dismissal_no_notice/yes/N.A2}; it ends mid-sentence and is counted as incorrect. Finish reason was not saved, so truncation cause is unverified. All other test answers parse. Stability can coexist with consistently wrong answers; it is not an accuracy measure.

\subsection*{A.5 Breakdown by intervention, direction, tier and schema}

\begingroup\small\setlength{\tabcolsep}{4pt}
\begin{longtable}{llrrrrrr}
\caption{Correct targets by intervention category; counts, not independent samples.}\\
\toprule
Model & Prompt & F & N3 & S & N & I1 & I2 \\
\midrule\endfirsthead
\toprule
Model & Prompt & F & N3 & S & N & I1 & I2 \\
\midrule\endhead
\bottomrule\endfoot
Sol & pA & 26/32 & 25/32 & 10/16 & 27/39 & 38/48 & 38/48 \\
Luna & pA & 27/32 & 26/32 & 12/16 & 31/39 & 40/48 & 42/48 \\
Haiku 4.5 & pA & 28/32 & 28/32 & 14/16 & 32/39 & 43/48 & 43/48 \\
Sol & pB & 27/32 & 25/32 & 12/16 & 27/39 & 39/48 & 39/48 \\
Luna & pB & 27/32 & 26/32 & 12/16 & 29/39 & 40/48 & 40/48 \\
Haiku 4.5 & pB & 27/32 & 28/32 & 14/16 & 35/39 & 44/48 & 45/48 \\
\end{longtable}\endgroup


\begingroup\small\setlength{\tabcolsep}{4pt}
\begin{longtable}{lllr}
\caption{M7 expanded into actual base-gold to target-gold directions.}\\
\toprule
Model & Prompt & Direction & Correct targets \\
\midrule\endfirsthead
\toprule
Model & Prompt & Direction & Correct targets \\
\midrule\endhead
\bottomrule\endfoot
Sol & pA & yes $\to$ no & 25/30 \\
Sol & pA & no $\to$ yes & 11/16 \\
Sol & pA & no $\to$ undetermined & 2/2 \\
Sol & pA & undetermined $\to$ yes & 4/10 \\
Sol & pA & undetermined $\to$ no & 19/22 \\
Luna & pA & yes $\to$ no & 27/30 \\
Luna & pA & no $\to$ yes & 10/16 \\
Luna & pA & no $\to$ undetermined & 2/2 \\
Luna & pA & undetermined $\to$ yes & 7/10 \\
Luna & pA & undetermined $\to$ no & 19/22 \\
Haiku 4.5 & pA & yes $\to$ no & 27/30 \\
Haiku 4.5 & pA & no $\to$ yes & 14/16 \\
Haiku 4.5 & pA & no $\to$ undetermined & 2/2 \\
Haiku 4.5 & pA & undetermined $\to$ yes & 9/10 \\
Haiku 4.5 & pA & undetermined $\to$ no & 18/22 \\
Sol & pB & yes $\to$ no & 26/30 \\
Sol & pB & no $\to$ yes & 11/16 \\
Sol & pB & no $\to$ undetermined & 2/2 \\
Sol & pB & undetermined $\to$ yes & 5/10 \\
Sol & pB & undetermined $\to$ no & 20/22 \\
Luna & pB & yes $\to$ no & 26/30 \\
Luna & pB & no $\to$ yes & 12/16 \\
Luna & pB & no $\to$ undetermined & 2/2 \\
Luna & pB & undetermined $\to$ yes & 5/10 \\
Luna & pB & undetermined $\to$ no & 20/22 \\
Haiku 4.5 & pB & yes $\to$ no & 27/30 \\
Haiku 4.5 & pB & no $\to$ yes & 12/16 \\
Haiku 4.5 & pB & no $\to$ undetermined & 2/2 \\
Haiku 4.5 & pB & undetermined $\to$ yes & 9/10 \\
Haiku 4.5 & pB & undetermined $\to$ no & 19/22 \\
\end{longtable}\endgroup


\begingroup\small\setlength{\tabcolsep}{4pt}
\begin{longtable}{lllrr}
\caption{Tier breakdown. Tiers are intended structural difficulty, not empirically validated levels.}\\
\toprule
Model & Prompt & Tier & M1 & M2 \\
\midrule\endfirsthead
\toprule
Model & Prompt & Tier & M1 & M2 \\
\midrule\endhead
\bottomrule\endfoot
Sol & pA & T1 & 9/12 (75.0\%) & 3/4 (75.0\%) \\
Sol & pA & T2 & 27/30 (90.0\%) & 14/15 (93.3\%) \\
Sol & pA & T3 & 18/28 (64.3\%) & 7/15 (46.7\%) \\
Sol & pA & T4 & 7/10 (70.0\%) & 3/5 (60.0\%) \\
Luna & pA & T1 & 12/12 (100.0\%) & 4/4 (100.0\%) \\
Luna & pA & T2 & 28/30 (93.3\%) & 14/15 (93.3\%) \\
Luna & pA & T3 & 18/28 (64.3\%) & 9/15 (60.0\%) \\
Luna & pA & T4 & 7/10 (70.0\%) & 4/5 (80.0\%) \\
Haiku 4.5 & pA & T1 & 12/12 (100.0\%) & 4/4 (100.0\%) \\
Haiku 4.5 & pA & T2 & 25/30 (83.3\%) & 14/15 (93.3\%) \\
Haiku 4.5 & pA & T3 & 23/28 (82.1\%) & 9/15 (60.0\%) \\
Haiku 4.5 & pA & T4 & 10/10 (100.0\%) & 5/5 (100.0\%) \\
Sol & pB & T1 & 11/12 (91.7\%) & 3/4 (75.0\%) \\
Sol & pB & T2 & 28/30 (93.3\%) & 14/15 (93.3\%) \\
Sol & pB & T3 & 18/28 (64.3\%) & 7/15 (46.7\%) \\
Sol & pB & T4 & 7/10 (70.0\%) & 3/5 (60.0\%) \\
Luna & pB & T1 & 11/12 (91.7\%) & 3/4 (75.0\%) \\
Luna & pB & T2 & 28/30 (93.3\%) & 13/15 (86.7\%) \\
Luna & pB & T3 & 19/28 (67.9\%) & 10/15 (66.7\%) \\
Luna & pB & T4 & 7/10 (70.0\%) & 3/5 (60.0\%) \\
Haiku 4.5 & pB & T1 & 12/12 (100.0\%) & 4/4 (100.0\%) \\
Haiku 4.5 & pB & T2 & 25/30 (83.3\%) & 15/15 (100.0\%) \\
Haiku 4.5 & pB & T3 & 24/28 (85.7\%) & 12/15 (80.0\%) \\
Haiku 4.5 & pB & T4 & 8/10 (80.0\%) & 4/5 (80.0\%) \\
\end{longtable}\endgroup


\begingroup\small\setlength{\tabcolsep}{4pt}
\begin{longtable}{lrrrrrr}
\caption{Test accuracy by schema, including base and all interventions.}\\
\toprule
Schema & Sol A & Luna A & Haiku A & Sol B & Luna B & Haiku B \\
\midrule\endfirsthead
\toprule
Schema & Sol A & Luna A & Haiku A & Sol B & Luna B & Haiku B \\
\midrule\endhead
\bottomrule\endfoot
refund\_window & 23/23 & 23/23 & 23/23 & 23/23 & 23/23 & 23/23 \\
witnessed\_contract & 18/23 & 22/23 & 19/23 & 21/23 & 21/23 & 23/23 \\
booking\_cancel & 28/30 & 30/30 & 30/30 & 30/30 & 29/30 & 28/30 \\
minor\_cancel & 30/30 & 30/30 & 28/30 & 30/30 & 30/30 & 30/30 \\
debt\_limitation & 22/30 & 24/30 & 26/30 & 22/30 & 22/30 & 27/30 \\
contract\_breach & 14/37 & 20/37 & 20/37 & 16/37 & 23/37 & 26/37 \\
dismissal\_no\_notice & 26/36 & 26/36 & 33/36 & 26/36 & 26/36 & 34/36 \\
hotel\_safe & 20/30 & 23/30 & 30/30 & 20/30 & 20/30 & 25/30 \\
\end{longtable}\endgroup


\begingroup\small\setlength{\tabcolsep}{4pt}
\begin{longtable}{llrrrrrrr}
\caption{Disagreement counts: expected gold versus observed answer (U = undetermined).}\\
\toprule
Model & Prompt & Y$\to$U & N$\to$U & Y$\to$N & N$\to$Y & U$\to$Y & U$\to$N & Invalid \\
\midrule\endfirsthead
\toprule
Model & Prompt & Y$\to$U & N$\to$U & Y$\to$N & N$\to$Y & U$\to$Y & U$\to$N & Invalid \\
\midrule\endhead
\bottomrule\endfoot
Sol & pA & 40 & 18 & 0 & 0 & 0 & 0 & 0 \\
Luna & pA & 22 & 15 & 4 & 0 & 0 & 0 & 0 \\
Haiku 4.5 & pA & 9 & 14 & 4 & 2 & 0 & 0 & 1 \\
Sol & pB & 36 & 15 & 0 & 0 & 0 & 0 & 0 \\
Luna & pB & 30 & 15 & 0 & 0 & 0 & 0 & 0 \\
Haiku 4.5 & pB & 12 & 2 & 4 & 5 & 0 & 0 & 0 \\
\end{longtable}\endgroup


Every Sol pA disagreement with gold (58/239) is a yes/no gold paired with an undetermined answer. This is a diagnostic pattern, not proof that all 58 answers are correct: the natural-language entailment must be adjudicated item by item.

\subsection*{A.6 Baselines and paired comparisons}

\begingroup\small\setlength{\tabcolsep}{4pt}
\begin{longtable}{lrrrrr}
\caption{Model-free test baselines. Oracle is a pipeline check, not a competing baseline.}\\
\toprule
Baseline & M0 & M1 & M2 & M3a & M3b \\
\midrule\endfirsthead
\toprule
Baseline & M0 & M1 & M2 & M3a & M3b \\
\midrule\endhead
\bottomrule\endfoot
const-yes & 8/24 (33.3\%) & 26/80 (32.5\%) & 16/39 (41.0\%) & 16/48 (33.3\%) & 16/48 (33.3\%) \\
const-no & 8/24 (33.3\%) & 52/80 (65.0\%) & 15/39 (38.5\%) & 16/48 (33.3\%) & 16/48 (33.3\%) \\
const-undetermined & 8/24 (33.3\%) & 2/80 (2.5\%) & 8/39 (20.5\%) & 16/48 (33.3\%) & 16/48 (33.3\%) \\
keyword-trigger & 24/24 (100.0\%) & 62/80 (77.5\%) & 6/39 (15.4\%) & 18/48 (37.5\%) & 18/48 (37.5\%) \\
oracle & 24/24 (100.0\%) & 80/80 (100.0\%) & 39/39 (100.0\%) & 48/48 (100.0\%) & 48/48 (100.0\%) \\
\end{longtable}\endgroup


The keyword-trigger baseline is deliberately gold-assisted: it starts from base gold and changes the answer when the edited sentence overlaps rule vocabulary. Its M0=100\% is by construction, not learned performance. It fails near-miss specificity. The criterion that a baseline should not solve all headline metrics excludes oracle, which must score 100\% by construction.

\begingroup\small\setlength{\tabcolsep}{4pt}
\begin{longtable}{lrrr}
\caption{Exploratory paired accuracy difference across all non-base items, percentage points [95\% CI].}\\
\toprule
Comparison & pA & pB & Pooled prompts \\
\midrule\endfirsthead
\toprule
Comparison & pA & pB & Pooled prompts \\
\midrule\endhead
\bottomrule\endfoot
Sol minus Haiku 4.5 & -11.2 [-19.1, -2.6] & -11.2 [-18.7, -2.5] & -11.2 [-18.0, -3.2] \\
Sol minus Luna & -6.5 [-11.1, -2.3] & -2.3 [-7.9, 1.4] & -4.4 [-9.2, -0.7] \\
Luna minus Haiku 4.5 & -4.7 [-11.6, 2.6] & -8.8 [-14.5, -2.8] & -6.7 [-12.6, -0.7] \\
\end{longtable}\endgroup


\begingroup\small\setlength{\tabcolsep}{4pt}
\begin{longtable}{lrrr}
\caption{Exploratory paired differences for each headline metric, pooled prompts.}\\
\toprule
Comparison & M1 & M2 & M3a \\
\midrule\endfirsthead
\toprule
Comparison & M1 & M2 & M3a \\
\midrule\endhead
\bottomrule\endfoot
Haiku 4.5 minus Sol & 8.8 [-2.0, 18.4] & 16.7 [3.4, 29.4] & 10.4 [-2.1, 21.9] \\
Luna minus Sol & 3.1 [0.5, 8.1] & 7.7 [-1.6, 19.8] & 3.1 [-2.1, 8.3] \\
\end{longtable}\endgroup


Intervals excluding zero describe differences in agreement with current gold only. They do not rescue gold validity, eliminate multiple-testing effects, or establish one model's superiority. The dev/test ordering reversal must be reported, not optimized away.

\subsection*{A.7 Archived development iterations and validity findings}

\begingroup\small\setlength{\tabcolsep}{4pt}
\begin{longtable}{llrr}
\caption{Saved superseded run snapshots. They overlap later runs; no comparable re-scoring without the exact historical dataset.}\\
\toprule
Snapshot & Model & Records & Invalid under current parser \\
\midrule\endfirsthead
\toprule
Snapshot & Model & Records & Invalid under current parser \\
\midrule\endhead
\bottomrule\endfoot
pre-goods-repair c & Haiku 4.5 & 235 & 0 \\
pre-goods-repair c & Luna & 235 & 0 \\
pre-goods-repair c & Sol & 235 & 0 \\
pre-expansion a & Haiku 4.5 & 119 & 0 \\
pre-expansion b & Haiku 4.5 & 119 & 0 \\
pre-expansion a & Luna & 119 & 0 \\
pre-expansion b & Luna & 119 & 0 \\
pre-expansion a & Sol & 119 & 0 \\
pre-expansion b & Sol & 119 & 0 \\
\end{longtable}\endgroup


Some intermediate 119-item and 235-item states were not separately snapshotted. The exact historical facts/gold are not versioned for every development round; re-scoring archived answers against today's repaired templates would be misleading. Therefore this table reports verifiable record/parse counts, not reconstructed performance. Small gateway verbatim/style probes were exploratory checks and were not a validated evaluation or fully archived experiment.

\begingroup\small\setlength{\tabcolsep}{4pt}
\begin{longtable}{p{3.4cm}p{4.7cm}p{7.1cm}}
\caption{Post-run audit: examples requiring semantic adjudication.}\\
\toprule
Schema & Gold assumption & What the English facts fail to establish \\
\midrule\endfirsthead
\toprule
Schema & Gold assumption & What the English facts fail to establish \\
\midrule\endhead
\bottomrule\endfoot
hotel\_safe & Watch in safe at theft & Deposited on first day does not entail remained in safe at theft. \\
debt\_limitation & No written acknowledgment & Never wrote to creditor does not exclude acknowledgment to someone else. \\
dismissal\_no\_notice & No written job protection & No extra agreement does not exclude protection inside the standard employment contract. \\
contract\_breach & Foreseeable/unforeseeable event & Annual recurrence or historical novelty does not logically settle foreseeability of the specific event. \\
Duplicate stimuli & 239 independent texts & There are 183 unique prompts; 46 duplicate groups contain 102 item IDs (56 excess copies). \\
M5 proxy & Edited fact is the unique deciding fact & Multiple supporting facts and output-format differences are not represented by the single-ID proxy. \\
\end{longtable}\endgroup


Conditional rates select different base-correct subsets across models, and ambiguous sentences can occur in multiple duplicate stimuli. Cluster resampling handles within-schema dependence but does not remove repeated-stimulus weighting or repair semantic errors. Some v1 flips also drop dependent facts, so not every generated pair is a one-proposition edit. Full detail is in \path{docs/TEST_VALIDITY_AUDIT_v1.md}. No items were removed after seeing test scores.

\subsection*{A.8 Revised implementation status and next experiment}

\begingroup\small\setlength{\tabcolsep}{4pt}
\begin{longtable}{lrrr}
\caption{V2 offline implementation only: no model evaluation performed.}\\
\toprule
Split & Schemas & Unique stimuli & Intervention edges \\
\midrule\endfirsthead
\toprule
Split & Schemas & Unique stimuli & Intervention edges \\
\midrule\endhead
\bottomrule\endfoot
Repaired inherited dev & 15 & 346 & 361 \\
New held-out test & 4 & 117 & 117 \\
Total & 19 & 463 & 478 \\
\end{longtable}\endgroup


The revised offline implementation now passes 14 regression tests. It separates stimulus IDs from intervention edges, checks exactly one fact-ID change per edge, and cross-checks all read-once three-valued formulas using an independent Boolean-completion evaluator. The v1 schemas are all development data in v2; four new schemas are reserved for a later evaluation. V2 retains the original insurance reporting-possibility exception with explicit possible/impossible facts, rather than silently changing it to another rule. Its runtime can retain full request/response metadata and distinguish prompt variation from repeat agreement. M5 is excluded from ability conclusions.

\textbf{Remaining work.} Offline checks do not establish natural-language or legal validity. Review the new items, resolve outstanding semantic issues with a second reader where possible, finalize the research question with the supervisor, and only then authorize the next API evaluation. The API gate remains disabled; no v2 model scores are claimed. This project has implemented a fixed-template diagnostic benchmark, not a reliable automated legal Challenger, and has not established novelty or superiority over existing faithfulness methods.

\end{document}
